\chapter{Upgrading, uninstalling, and the licences you ship}\label{ch:upgrades}

\section{Upgrade order}

Upgrades install in place and keep the config and all the data. The order that has to be followed
is:

\begin{enumerate}
  \item \textbf{Back up the database and the config} (chapter~\ref{ch:database}). An upgrade never
    touches them, but a migration that turns out to be the wrong one has no built-in undo.
  \item \textbf{Install the new artifact over the old one.} The Windows installer, \texttt{install.sh}
    and the \texttt{.deb} postinst each run \texttt{ate setup} as their last step, and that runs the
    migration, so an install through those three leaves the database current. The database it
    migrates is \texttt{\textless data\textgreater/ate.db}, which is what the default config points
    at; a config with a custom \texttt{database.path} is not migrated, so a custom database still
    needs \texttt{ate db migrate} after the install. Step 4 below is what catches it. A wheel
    installed with \texttt{pip} or \texttt{pipx} runs nothing: its database stays at the old schema
    until you migrate it, and ATE refuses to start against a database older than itself.
  \item \textbf{Migrate, when the artifact did not, or when \texttt{database.path} points elsewhere.}
    \texttt{ate db migrate} is idempotent: run it on an already-current database and it succeeds,
    prints the version, and changes nothing.
  \item \textbf{Check with \texttt{ate db status}.} It reads the configured path, so it also catches a
    custom database the installer never touched. On this build it prints \texttt{5} and exits 0 ---
    \texttt{5} is the schema version, not the ATE version, so do not compare the two. A missing or
    older database exits 1 with the refusal to run, which is the migration not having happened.
  \item \textbf{Check the Media Releases.} A Media Release whose files no longer match its recorded
    checksums warns rather than failing; after updating a release, run \texttt{ate definition update}
    on its \texttt{.atedef} (chapter~\ref{ch:media-releases}).
\end{enumerate}

Each artifact clears its own program files before writing the new ones, so nothing stale survives an
upgrade: the Windows installer deletes \texttt{\_internal} first; \texttt{install.sh} removes the
bundle's \texttt{ate} and \texttt{\_internal} before copying; dpkg removes everything the package
installed. Config and data live outside the install directory on every platform, so none of this
touches them.

Going backwards is not supported by the database. A schema version is only ever increased, and an
older ATE started against a database built by a newer one refuses with \texttt{this ATE is too old
for the database at \dots; update ATE}. If you must roll back an ATE version, roll the database back
with the backup you took.

The version ATE reports is the git tag the bundle was built from, in semver form, stamped in at
build time; prerelease tags are written the PEP 440 way (\texttt{v0.1.0-rc1} reports
\texttt{0.1.0rc1}, and the \texttt{.deb} names that \texttt{0.1.0\textasciitilde rc1}). That version
is recorded on every Test Run and printed in the Test Report footer, so a record says which ATE
produced it.

\section{Uninstalling}

Uninstall removes the program, and only the program.

\begin{itemize}
  \item Windows: the uninstaller deletes the install directory and takes its own entry back out of
    the machine PATH. It leaves any PATH entry it did not add, and writes the value back as
    \texttt{REG\_EXPAND\_SZ} so entries like \texttt{\%USERPROFILE\%} keep expanding.
  \item \texttt{apt remove}: config and data stay. \texttt{apt purge}: removes
    \texttt{/etc/ate/config.json} and the config directory if it is then empty, and never touches
    \texttt{/var/lib/ate}.
  \item \texttt{./install.sh --uninstall --prefix /opt/ate}: removes the bundle and its link. It
    refuses to remove a prefix that has no ATE executable in it, so it cannot wipe a directory an
    operator pointed it at by mistake.
  \item The wheel: \texttt{pipx uninstall ate}. Whatever \texttt{ate setup} created is not part of
    the wheel and stays behind.
\end{itemize}

Test Runs are records. No uninstall path deletes the database, its WAL files, the reports or the
logs, and if you want them gone you remove them yourself, having checked whether anyone needs the
records.

\section{Third-party licences, and where they land}

ATE's interface and its Test Report renderer are Qt, used through PySide6, both LGPL. Whoever
receives a Station needs to be able to read that licence, see which Qt is in the bundle, find its
source, and know that Qt can be replaced. Four files carry that, generated from the wheels actually
installed in the build rather than written by hand:

\begin{lstlisting}[language=ATEShell]
python packaging/licence/notices.py --out build/licences
\end{lstlisting}

\texttt{NOTICE-Qt-PySide6.txt} (what Qt, PySide6 and shiboken6 versions are in the bundle, their
licence as the wheels declare it, where the source is, and the statement that the shared libraries
are copied unmodified and can be replaced by another build of the same Qt version);
\texttt{THIRD\_PARTY\_NOTICES.txt} (every runtime dependency, its version and the licence its wheel
names, walking every requirement of every requirement); \texttt{LGPL-3.0.txt} and
\texttt{GPL-3.0.txt}, the full texts. A licence a wheel does not name, or that the script cannot
name, is reported as such rather than guessed --- the file ends with a list of the ones to check
upstream, and the Lua that lupa embeds is named with its own licence link.

In the bundle they sit under \texttt{\_internal/licences/}, and \texttt{verify\_bundle} refuses to
release a bundle that is missing any of the four. The tarball and the \texttt{.deb} ship that
directory, and the suite checks both artifacts still carry it. The wheel is not checked for them by
the suite: it is a Python package whose dependencies come from PyPI, and nothing here proves that
the notices travel with it. Say so if compliance depends on it.

\cando{upgrade a Station in the right order, uninstall it without taking a record with it, and say
where the licence texts a Station receives are to be found.}
